\section{Introduction}

When we constructed the first directed citation graph of the huashen218 bidirectional alignment corpus, we discovered that the publicly available reading list contained 49 extractable paper identifiers --- not the $\sim$400 from the underlying systematic review.
What we could measure, however, told a notable story: every HCI alignment paper in our resolved set cited ML/NLP work as its sole within-corpus outgoing reference, while ML/NLP alignment papers directed only 10.6\% of their outgoing citations toward HCI.
A directed citation ratio of 0.107.
The infrastructure to detect this asymmetry is now validated. The full statistical test awaits a larger corpus.

The framing of ``bidirectional AI alignment'' --- the idea that ML/NLP and HCI communities should mutually adapt to one another --- rests on an implicit assumption: that these communities are intellectually coupled.
Cross-community citation behavior is a structural indicator of epistemic integration.
If alignment researchers from ML/NLP and HCI communities genuinely engage with each other's work, their citation networks should reflect symmetric knowledge flow.
If citation is asymmetric --- HCI researchers consistently citing ML/NLP but not vice versa --- then bidirectionality may be more aspirational than structural.

This question is newly tractable. \citet{Shen2024Position} assembled the huashen218 corpus: a curated systematic bibliography of 400+ interdisciplinary alignment papers spanning NeurIPS, ICML, CHI, CSCW, and adjacent venues.
The corpus operationalizes the ``bidirectional alignment community'' as a named set of papers.
Measuring directed citation behavior \emph{within} this corpus is the first empirical test of whether the community is citation-integrated.

The measurement, however, exposes a deeper problem.
Standard bibliometric approaches --- classifying papers by venue name substring matching --- achieve only \textbf{12\% label coverage} on this corpus.
The reason is mundane but consequential: the Semantic Scholar Academic Graph (S2AG), the public API for directed citation data, returns full proceedings names (``Proceedings of the 36th Annual Conference on Neural Information Processing Systems'') rather than abbreviated venue strings (``NeurIPS'').
Substring matching silently fails on the long-form variants, leaving 88\% of papers unclassified and the measurement impossible.

This failure points to the gap: there is no validated pipeline or classification method for constructing directed citation graphs within interdisciplinary alignment corpora.
Venue-string approaches, the default in most bibliometric tools, are not suitable here.
The classification infrastructure must be built before any citation asymmetry can be measured.

\textbf{Our key insight:} Switching the primary classifier from venue name strings to S2AG's \texttt{fieldsOfStudy} tags --- which encode disciplinary assignment at the knowledge-graph level, independent of venue name formatting --- achieves \textbf{100\% paper coverage} on the same corpus, compared to 12\% for venue strings. This is a three-line change with an eight-fold coverage improvement.

Building on this insight, we make the following contributions:

\begin{enumerate}
\item \textbf{A validated S2AG bibliometric pipeline} for within-corpus directed citation analysis: batch ID resolution via \texttt{POST /paper/batch}, FoS-primary classification (Scheme 3), within-corpus DiGraph construction using NetworkX, and gate-controlled evaluation. The pipeline passes 23/23 unit tests and is fully cached for zero-API re-execution.

\item \textbf{FoS-primary classification (Scheme 3)}: a method that uses S2AG \texttt{fieldsOfStudy[0]} as the primary venue-group classifier, achieving 100\% label coverage on the huashen218 corpus vs.\ 12\% for venue-string-only approaches. The method is domain-general and applicable to any interdisciplinary corpus indexed by S2AG.

\item \textbf{The first directional citation measurement within the huashen218 alignment corpus}: in the 33-paper resolved proxy, ML/NLP$\to$HCI proportion = 10.6\%, HCI$\to$ML/NLP proportion = 100\%, ratio $\approx$ 0.107 --- directionally consistent with asymmetric epistemic dependency. No statistical falsifier was triggered; full statistical testing requires the augmented corpus (h-e1-v2).

\item \textbf{A methodological finding on corpus proxy construction}: the huashen218 GitHub reading list yields 49 extractable identifiers from $\sim$130 linked papers --- 12.5\% of the target corpus. Public reading lists from systematic reviews are insufficient bibliometric dataset proxies; programmatic S2AG citation-of-citation search is the viable reconstruction approach.
\end{enumerate}

We organize the paper as follows.
Section~\ref{sec:related} situates our work in the context of prior cross-community citation analysis and bibliometric infrastructure.
Section~\ref{sec:method} describes the pipeline design and the classification schemes.
Section~\ref{sec:experiments} presents experimental setup.
Section~\ref{sec:results} reports results.
Section~\ref{sec:discussion} discusses implications and limitations.
Section~\ref{sec:conclusion} concludes.
